\chapter{Метод постквантово-посиленої OPAQUE-автентифікації}
\label{ch:opaque}

% ============================================================
%  РОЗДІЛ 5 -- криптографічний контур, частина 1: автентифікація.
%  Джерело: aura-security-opaque-rs + стаття main_jisa.tex.
%  Найбільш підготовлений до публікаційного використання матеріал (готовність 88/100).
% ============================================================

\section{Постановка задачі та модель загроз автентифікації}
\label{sec:opaque-threat}
% -- Чому середовище виконання політик не можна будувати на слабкій автентифікації.
% -- Класичний і квантовий супротивник; компрометація БД сервера.

\section{Архітектура трьохповідомленнєвого протоколу}
\label{sec:opaque-arch}
% -- Гібридний PQ-OPAQUE: три повідомлення; нормативний формат передавання і правила відхилення.

\section{Реєстрація та зв'язування ідентифікатора й контексту}
\label{sec:opaque-registration}
% -- Фаза реєстрації; зв'язування ідентифікатора облікового запису з контекстом.

\section{Гібридне DH/KEM-узгодження сесії}
\label{sec:opaque-hybrid}
% -- X25519 + ефемерний ML-KEM-768 [cite:fips203]; перевірка зв'язування протоколу з протоколограмою.

\section{Операційна модель \texttt{oprf\_seed} і межі квантового впливу}
\label{sec:opaque-oprf}
% -- Чітке операційне окреслення \texttt{oprf\_seed}; що саме захищає; межі постквантового підсилення.
% -- Argon2id як рівень ускладнення перебору пароля [cite:rfc9106].

\section{Формальний аналіз і межі доведення}
\label{sec:opaque-formal}
% -- Tamarin [cite:tamarin]; перелік розряджених властивостей; межі доведення.

\section{Реалізація, простежуваність специфікації та вимірювання швидкодії}
\label{sec:opaque-impl}
% -- Реалізація на Rust; простежуваність від специфікації до коду;
%    безпека пам'яті, обнулення секретів; розміри повідомлень, затримка.
Метод супроводжується еталонною реалізацією на Rust із простежуваністю специфікації до коду, перевіркою безпеки пам'яті, обнуленням секретів та відтворюваними вимірюваннями швидкодії примітивів і протоколу. \emph{[Розкрити кількісні результати.]}

\rozdilconcl{5}
\emph{[Висновки до розділу 5.]}
